\begin{abstract}
山区洪涝灾害会同时破坏道路通行条件和地面通信覆盖，使应急物资投送从单纯的最短路问题转化为地形、能量、时限、装备周转和连续通信共同约束的协同决策问题。题目给出15个服务区、80个不可拆货箱、三类运输无人机及共享电池，并允许在直连通信不足时使用中继无人机。本文按照“单点能力标定—多点运输调度—运输通信协同—独立任务分区”的递进关系，建立统一的航段、能耗和资源事件模型。模型始终以硬约束可行为第一层目标，再比较一般物资及时性、任务完成时间、能耗和架次数。

针对问题一，首先沿每条航段遍历实际经过的30~m DEM像元，以最高地形加50~m确定巡航海拔；结合载荷相关等效航程、爬升能耗和返航安全余量，利用二分法求解“服务区—机型”最大安全载荷。随后枚举满足质量、体积和能量约束的完整货箱子集，以架次数、能耗和累计作业时间为字典序代价，通过集合划分动态规划获得单点组批方案。结果表明，45个“服务区—机型”组合均可空载返航，但B型和C型在部分远距或高地形节点受到能量限制；全体物资共需18架次，累计能耗59.240~kWh，累计作业时间9.111~h。18个批次中，体积、质量和能量分别成为9、5和4个批次的最紧约束，说明不可拆货箱条件下仅按额定载质量估计架次数会产生偏差。

针对问题二，将货箱组批、访问顺序、机型、实体机、共享电池和开始时刻置于同一事件时间轴。逐站递推交付时刻与剩余载荷，医疗及首批保障货箱采用硬截止约束，一般物资采用归一化加权迟到指标；实体机占用至返航，电池则占用至两阶段充电结束。求解时分别构造硬箱专送基线与同目的地混装方案，通过软箱补入、单区精确组批和双区串访形成两套候选，并由事件仿真器完成资源排程。与26架次硬箱专送基线相比，所选混装方案以22架次完成80箱唯一交付，31个硬时限箱均无违约，一般物资超期箱由23箱降至9箱，$WTD$由111.492降至38.787；完成时间为2.679~h，总能耗为69.602~kWh，分别较基线下降24.1\%和16.2\%。

针对问题三，依据双向链路预算、自由空间损耗和DEM视线遮挡建立直连与单中继可用判据，将运输轨迹展开为爬升、巡航、下降和交接区间。中继悬停点须同时满足边界、回传链路、往返能量和300~m离地高度要求；中继任务计入准备、飞抵、建链、悬停通信、返航、周转和能源组件充电。本文以问题二路线为已验证候选，采用离线筛选后固化的六个点位和七个波次组织中继资源，并用连续区间证书进行全时段认证，0.5~s网格仅作为独立复核。最终22个运输架次被组织为7个波次，配置11个中继架次，由2架中继机和5套能源组件执行；68个通信分段包含41段直连和27段中继，连续认证无未解析区间，65532个加密采样点无中断。联合完成时间为6.572~h，总能耗79.097~kWh，最小保守链路余量为0.1397~dB。该结果证明当前方案可行，但不构成运输—通信全局最优性证明。

针对问题四，冻结问题三的组批、路线、时序和实际通信指派，以同架次共访关系构造服务区图，得到11个不可拆原子块；对每个合法分区，从通信分段反推组内中继任务，并以固定占用区间的峰值并发数计算各类资源最小配置。实际穷举1023个两组候选和57002个三组带标签候选，后者对应28501个无标签等价类。结果表明：两组最优方案不产生库存缺口，三组最优方案需要额外增加1架中继无人机，并相对集中执行多配置1套中继能源组件。因此，在“先消除库存缺口，再减少总配置并改善均衡”的评价准则下选择两组方案；若强调区域自治，则需接受额外资源。

模型从航段像元、逐架次能量、80箱唯一覆盖、资源无冲突和连续通信五个层次进行检验。敏感性实验表明：返航余量提高至25\%会使问题一架次数由18增至19；既定问题二方案在1200~s事件延误内仍满足硬时限；问题三原证书可承受0.1~dB附加损耗。本文以统一事件时间轴贯通运输、能源和通信，使结论能够回溯到具体货箱、架次或时间区间。

\keywords{应急物资运输；集合划分动态规划；共享电池调度；连续通信认证；任务分区}
\end{abstract}
